\section{Limitations}
\label{sec:limitations}

\paragraph{The working posterior is not identified by a probability map.}
The shared-threshold construction selects one coupling from many with the same pixel
marginals. Marginal calibration, even if exact, does not identify expected nonlinear
geometric loss. Repeated annotations or another source of label variability would be
needed to validate this coupling directly; one reference mask per image cannot
separate marginal calibration error from dependence-model error.

\paragraph{The loss convention is substantive.}
Dice is naturally bounded, but surface distances require an explicit empty-mask
rule. Assigning a one-sided empty comparison the image diagonal, a fixed cap, or an
application-specific cost produces different losses and can change the ranking. Our
normalized penalized HD95 is one predeclared choice. The published SDC baseline's
empty--empty Dice value also differs from ours, so its theoretical approximation
guarantee does not automatically cover empty deployed masks under our loss.
Dropping undefined cases would change the target population and is not equivalent.

\paragraph{Quadrature control is not posterior or AURC control.}
Equation~\eqref{eq:decomp} separates numerical and posterior terms but does not make
the latter small by assertion. The variation constant can also be loose. A larger
\(M\) improves the worst-case midpoint approximation rate to the \(Q_p\) integral;
it need not improve empirical AURC monotonically. Turning value error into an AURC
guarantee additionally requires rank-stability or error-tail control; a pairwise
margin is one sufficient route for exact ordering, not a consequence of the bound.

\paragraph{Scope and information limits.}
We focus on native binary masks and single-map, post-hoc confidence. If two cases
produce the same \((p,\Yhat)\), every deterministic score studied here gives them
the same confidence even when external context makes one more likely to fail.
Features, ensembles, auxiliary failure predictors, distribution-shift detectors,
multiclass actions, and multi-rater posterior validation are outside the present
scope. Finally, AURC measures ranking; it does not imply that a confidence value is
calibrated as a probability or an absolute forecast of loss.

\paragraph{Empirical generality and training variability.}
Three binary datasets and two model families cannot establish that metric-aligned
confidence will dominate across architectures, shifts, or annotation protocols.
Our paired comparisons deliberately condition on each fixed probability map and
therefore isolate confidence ranking, not segmentation-model accuracy.  The
paired image-cluster bootstrap quantifies held-out image variation for a fixed
fitted checkpoint; its percentile intervals and tail probabilities are approximate,
and with one predeclared training seed they do not include optimization or
training-sample variability.  Repeated training runs and external-site validation
would be needed for that broader claim. Moreover, the 15,876 condition-specific
rows reuse 4,069 unique held-out images across model conditions; they are not
15,876 independent images.

\section{Conclusion}
\label{sec:conclusion}

Selective segmentation confidence is meaningful only relative to a deployed mask
and deployment loss. We defined the corresponding ideal confidence, represented a
level-set construction as an explicit working posterior, and obtained one quadrature
form for any bounded measurable segmentation loss. The analysis separates numerical
approximation from disagreement between the working and true label laws. Dice
illustrates why this separation matters: expected Dice under the level-set law and
SDC's ratio of expected components are different estimands with structurally
non-nested conditions for approximating the ideal score. Empirically, nHD95-M32
consistently improves on Dice-M32 for nHD95 risk, but their Dice-risk ordering
splits six to four and SDC is the strongest primary score in six conditions.
Loss indexing therefore provides a general construction and a useful diagnostic
axis, not a guarantee that a same-named confidence dominates every baseline or
cross-loss score.
